\documentclass[11pt,letterpaper]{amsart}
\usepackage[T1]{fontenc}
\usepackage{lmodern}
\usepackage[margin=1.1in]{geometry}
\usepackage{amsmath,amssymb,amsthm,mathtools}
\usepackage{microtype}
\usepackage{xcolor}
\usepackage{hyperref}
\hypersetup{colorlinks=true,linkcolor=blue!45!black,urlcolor=blue!45!black,citecolor=blue!45!black,
 pdftitle={Countable strongly Choquet--Deny groups are precisely the virtually FC groups},
 pdfauthor={}}
\newtheorem{theorem}{Theorem}[section]
\newtheorem{lemma}[theorem]{Lemma}
\newtheorem{proposition}[theorem]{Proposition}
\newtheorem{corollary}[theorem]{Corollary}
\theoremstyle{definition}
\newtheorem{definition}[theorem]{Definition}
\newtheorem{example}[theorem]{Example}
\theoremstyle{remark}
\newtheorem{remark}[theorem]{Remark}
\newcommand{\R}{\mathbb R}
\newcommand{\N}{\mathbb N}
\newcommand{\Z}{\mathbb Z}
\newcommand{\F}{\mathbb F}
\DeclareMathOperator{\co}{co}
\DeclareMathOperator{\Orb}{Orb}
\DeclareMathOperator{\Prob}{Prob}
\DeclareMathOperator{\Meas}{Meas}
\DeclareMathOperator{\ran}{ran}
\DeclareMathOperator{\Span}{span}
\DeclareMathOperator{\FC}{FC}
\DeclareMathOperator{\supp}{supp}
\title[The classification of countable strongly Choquet--Deny groups]{Countable strongly Choquet--Deny groups
are precisely the virtually FC groups}
\date{23 October 2026}
\begin{document}
\begin{abstract}
We show that the countable strongly Choquet--Deny groups are precisely
the virtually FC groups, extending the finitely generated classification
of Dantas and Doucha. The proof uses their criterion for invariant
strictly convex renormability and the analytic construction in their
Heisenberg example. The positive direction replaces commutativity by
finite conjugacy averaging. For the converse, we construct the sequence
required by their obstruction whenever the FC-center has infinite index.
The same characterization holds for their positive strongly
Choquet--Deny condition.
\end{abstract}
\maketitle

\section{Introduction}
Dantas and Doucha introduced strongly Choquet--Deny groups and proved
that a countable group has this property if and only if every uniformly
bounded representation on a separable real Banach space admits an
equivalent invariant strictly convex norm \cite[Theorem~4.2]{DD}.
Their work follows the study of invariant renormings in \cite{D};
the connection between harmonic vectors and strict convexity also
appears in Jaworski--Neufang \cite[Proposition~2.5]{JN}.
Dantas--Doucha classified the finitely generated strongly Choquet--Deny
groups as the virtually abelian groups \cite[Theorem~5.6]{DD}.
This note extends their classification to arbitrary countable groups.

Recall that a group is \emph{FC} if all its conjugacy classes are finite,
and \emph{virtually FC} if it has an FC subgroup of finite index.
Its FC-center $\FC(G)$ is the subgroup of elements with finite conjugacy
class in $G$. The strong Choquet--Deny conditions in the following
statement are those of \cite[Definition~2.4]{DD}, recalled in
Section~\ref{sec:prelim}.

\begin{theorem}\label{thm:classification}
For a countable group $G$, the following conditions are equivalent:
\begin{enumerate}
\item $G$ is virtually FC.
\item Every uniformly bounded representation of $G$ on a separable real
Banach space admits an equivalent $G$-invariant strictly convex norm.
\item $G$ is strongly Choquet--Deny.
\item $G$ is positively strongly Choquet--Deny.
\end{enumerate}
\end{theorem}

A finitely generated FC group has a finite-index center: intersect
the centralizers of a finite generating set. Since a finite-index
subgroup of a finitely generated group is finitely generated, virtual
FC and virtual abelianness coincide in that setting. They differ for
general countable groups; for example, $\bigoplus_{n\geq1}S_3$ is FC
but not virtually abelian (Example~\ref{ex:finite-sum}).

\subsection{Proof and relation to earlier work}
We use two extensions of the methods of \cite{DD}.
For the positive direction, finite conjugacy averages show that a
midpoint decomposition in an absolutely convex orbit hull preserves
all norms and linear relations in the cyclic representation. This
produces two contractions whose average is the identity, extending
the abelian operator argument of \cite[Theorem~5.3]{DD}. Kakutani's
extremality theorem gives orbit extremality, and the criterion in
\cite[Theorem~3.1]{DD} supplies the invariant strictly convex norm.
Finite-index extension then treats virtually FC groups.

For the converse, \cite[Example~5.5, Claim~2]{DD} constructs an
obstruction from a sequence satisfying the three product conditions
in Claim~1 of that example. Proposition~\ref{thm:certificate} records
this implication, with a minor change to the damping function.
The additional group-theoretic argument is
Theorem~\ref{thm:sequence}, which constructs such a sequence whenever
$[G:\FC(G)]=\infty$. It is a finite extension argument using
centralizer cosets and the equation-avoidance lemma
(Lemma~\ref{lem:avoid}). Its relation to the rewritability theorem of
Blyth--Rhemtulla \cite{BR} and the switching constructions of
\cite{FTVF,FHTVF,EK} is explained in Section~\ref{sec:sequence}.

For comparison, a countable group is \emph{Choquet--Deny} if every
bounded $\mu$-harmonic function is constant on left cosets of the
subgroup generated by $\supp(\mu)$, for every probability measure
$\mu$ \cite[Definition~2.2]{DD}. These are precisely the groups
with no nontrivial ICC quotient \cite[Theorem~1]{FHTVF}, or
equivalently the hyper-FC-central groups. Here ICC means nontrivial
with every nonidentity conjugacy class infinite; hyper-FC-central
means that iterating FC-centers in successive quotients, with unions
at limit stages, exhausts the group. The positive direction of this
ordinary classification goes back to \cite[Theorems~4.5 and~4.8]{J}.
Theorem~\ref{thm:classification} requires finite index already at the
first FC-center. The examples in Section~\ref{sec:examples} show that
this is a strictly stronger requirement.

The equivalence of the positive and signed strong conditions for the
same countable group is a consequence of the classification.
The earlier relation in \cite[Proposition~2.6]{DD} identifies the
signed condition for $G$ with the positive condition for
$G\times\Z_2$.

\section{Preliminaries}\label{sec:prelim}
All groups are discrete, and all Banach spaces and function spaces are
over $\R$. All convex hull closures are taken in norm. A representation
$\pi:G\to\mathrm{GL}(X)$ is \emph{uniformly bounded} if
$M=\sup_{g\in G}\|\pi(g)\|<\infty$. In this case
\[
 \|x\|_0=\sup_{g\in G}\|\pi(g)x\|
\]
is an equivalent invariant norm, with
$\|x\|\leq\|x\|_0\leq M\|x\|$. Thus it suffices to treat isometric
representations. A norm $N$ is \emph{strictly convex} if
$N((x+y)/2)<1$ whenever $x\ne y$ and $N(x)=N(y)=1$.

\subsection{Measures and translations}
We follow \cite[Section~2.2]{DD} in writing $\Prob(G)$ for the
probability measures and $\Meas(G)=B_{\ell^1(G)}$ for the real signed
measures of total variation at most one. For $\mu\in\ell^1(G)$ set
$\mu^*(h)=\mu(h^{-1})$. Their convolution convention is
\[
 (f*\mu)(g)=\sum_{h\in G}f(gh)\mu(h^{-1}),\qquad
 (f*\mu^*)(g)=\sum_{h\in G}\mu(h)f(gh).
\]
We also write $(R_hf)(g)=f(gh)$ for right translation. Thus
$R_hR_k=R_{hk}$ and $f*\mu^*=\sum_h\mu(h)R_hf$, with convergence
in supremum norm. In particular, right convolution preserves the
closed span of the right translates of $f$. A bounded function is
$\mu$-harmonic if $f*\mu^*=f$. The symbol $\delta_g$ denotes a point
mass or the indicator of $\{g\}$, according to the ambient space.

\begin{definition}[{\cite[Definition~2.4]{DD}}]\label{def:strong}
A countable group $G$ is \emph{strongly Choquet--Deny} if the following
implication holds. Let
$\mu_n,\nu_n,\nu_n'\in\Meas(G)$ with $\mu_n=(\nu_n+\nu_n')/2$,
and let $f\in\ell^\infty(G)$ be fixed. If
\begin{equation}\label{eq:hypotheses}
 \|f*\mu_n^*-f\|_\infty\longrightarrow0,
 \qquad
 \sum_{n=1}^\infty\|d_{n+1}-d_n\|_\infty<\infty,
 \quad d_n=f*(\nu_n^*-(\nu_n')^*),
\end{equation}
then $\|d_n\|_\infty\to0$. The group is \emph{positively strongly
Choquet--Deny} if the same implication holds with all measures in
$\Prob(G)$.
\end{definition}

The summability in \eqref{eq:hypotheses} makes $(d_n)$ converge in
supremum norm, so this is equivalent to the subsequence conclusion in
\cite[Definition~2.4]{DD}. Restricting the measures shows that the
signed condition implies the positive one \cite[Proposition~2.5]{DD}.

\subsection{Orbit hulls and the FC-center}
Write $\Orb_G(x)=\{\pi(g)x:g\in G\}$ and
$\Orb_G(x,-x)=\Orb_G(x)\cup\Orb_G(-x)$, as in
\cite[Section~2.1]{DD}. Its closed convex hull is the closed absolutely
convex hull of $\Orb_G(x)$. We use the following criterion.

\begin{theorem}[{Dantas--Doucha, \cite[Theorem~3.1]{DD}}]
\label{thm:criterion}
An isometric representation of a group $G$ on a separable real Banach
space $X$ admits an equivalent invariant strictly convex norm if and
only if every $x\in X$ is an extreme point of
$\overline{\co}(\Orb_G(x,-x))$.
\end{theorem}

For countable groups, \cite[Theorem~4.2]{DD} further identifies the
strongly Choquet--Deny property with the existence of such a norm in
every separable uniformly bounded representation. We use this
equivalence in the proof of Theorem~\ref{thm:classification}.

Necessity in Theorem~\ref{thm:criterion} holds without separability:
all points of this hull have
invariant norm at most that of $x$, so strict convexity excludes a
nontrivial midpoint at $x$. In our converse the endpoints will even
lie in $\overline{\co}(\Orb_G(x))$.

For $g\in G$ let $C_G(g)=\{h\in G:hg=gh\}$ be its centralizer.
The FC-center is
\[
 \FC(G)=\{g\in G:[G:C_G(g)]<\infty\}.
\]
The index equals the cardinality of the conjugacy class of $g$.
The FC-center is a characteristic subgroup: conjugates of a product
belong to the product of the conjugacy classes of its factors, and
inversion and automorphisms preserve finiteness of conjugacy classes.
See also \cite[Section~5.B]{EK}. Moreover,
\begin{equation}\label{eq:fc-index}
 G\text{ is virtually FC}\quad\Longleftrightarrow\quad
 [G:\FC(G)]<\infty.
\end{equation}
For the nontrivial implication, let $H\leq G$ be FC of finite index.
The $G$-conjugacy class of $h\in H$ is a finite union of conjugates
of its finite $H$-conjugacy class. Thus $H\subseteq\FC(G)$.

\section{Invariant renormings for virtually FC groups}\label{sec:positive}
We prove the positive direction of Theorem~\ref{thm:classification}.
No countability assumption on the group is needed in this section.
The proof follows the contractive-operator strategy of
\cite[Theorem~5.3]{DD}. Finite conjugacy averaging also appears in
Jaworski's unitary-representation argument \cite[proof of Theorem~4.4]{J}.
Here we average squared norms of group-algebra translates to preserve
all linear relations in a cyclic Banach representation.

\begin{lemma}[{Kakutani; see \cite[Theorem~8.28]{EFHN}}]\label{lem:identity}
For every real Banach space $E$, the identity is an extreme point of
the closed unit ball of $\mathcal B(E)$.
\end{lemma}
\begin{proof}
We use the geometric proof in \cite[Theorem~8.28]{EFHN}.
Suppose $A,B:E\to E$ are contractions with $(A+B)/2=I$.
For every extreme point $\varphi$ of the dual unit ball $B_{E^*}$,
\[
 \varphi=\frac{A^*\varphi+B^*\varphi}{2},\qquad
 A^*\varphi,B^*\varphi\in B_{E^*}.
\]
Extremality gives $A^*\varphi=B^*\varphi=\varphi$.
By the Kre\u{\i}n--Milman theorem, $B_{E^*}$ is the weak-star closed
convex hull of its extreme points. Since adjoint operators are
weak-star continuous, $A^*$ and $B^*$ fix the whole dual unit ball.
Thus $A=B=I$.
\end{proof}

\begin{theorem}[FC orbit extremality]\label{thm:fc}
Let $G$ be an FC group. In every real isometric Banach representation,
without a separability assumption, every vector $x$ is extreme in its
closed absolutely convex orbit hull $\overline{\co}(\Orb_G(x,-x))$. Consequently
every uniformly bounded representation of $G$ on a separable real
Banach space admits an equivalent invariant strictly convex norm.
\end{theorem}
\begin{proof}
Fix $x=(y+z)/2$ with $y,z\in \overline{\co}(\Orb_G(x,-x))$. For a finitely supported element
$a=\sum_g a_g g$ of the real group algebra, write
$\pi(a)=\sum_g a_g\pi(g)$ and let
\[
 \mathcal O(a)=\{h^{-1}ah:h\in G\}.
\]
This set is finite: the support of $a$ is finite, and every element of
that support has finitely many conjugates. Define
\[
 \Phi_a(v)=\frac1{|\mathcal O(a)|}
                  \sum_{b\in\mathcal O(a)}\|\pi(b)v\|^2.
\]
This function is continuous, even, convex, and $G$-invariant. Indeed,
$\pi(b)\pi(h)=\pi(h)\pi(h^{-1}bh)$, the first factor is an isometry,
and conjugation permutes $\mathcal O(a)$.

Since $y,z\in \overline{\co}(\Orb_G(x,-x))$, continuity, convexity, and invariance give
$\Phi_a(y),\Phi_a(z)\leq\Phi_a(x)$. Thus
\[
 \Phi_a(x)\leq\frac{\Phi_a(y)+\Phi_a(z)}2\leq\Phi_a(x).
\]
All the individual squared-norm Jensen gaps in this finite average are
nonnegative, so they all vanish. Indeed, the triangle inequality gives
$\|\pi(b)x\|\leq(\|\pi(b)y\|+\|\pi(b)z\|)/2$, and hence
for each $b\in\mathcal O(a)$,
\[
 \frac{\|\pi(b)y\|^2+\|\pi(b)z\|^2}{2}-\|\pi(b)x\|^2
 \geq\frac14\bigl(\|\pi(b)y\|-\|\pi(b)z\|\bigr)^2.
\]
In particular, because $a\in\mathcal O(a)$,
\begin{equation}\label{eq:all-a}
 \|\pi(a)y\|=\|\pi(a)z\|=\|\pi(a)x\|
 \quad\text{for every finitely supported real }a.
\end{equation}

Put $E=\overline{\Span}\{\pi(g)x:g\in G\}$. Both $y$ and $z$ lie
in $E$. On its dense algebraic cyclic subspace define
\[
 A(\pi(a)x)=\pi(a)y,\qquad B(\pi(a)x)=\pi(a)z.
\]
In particular, \eqref{eq:all-a} gives
$\pi(a)x=0\Rightarrow\pi(a)y=\pi(a)z=0$. Applying this to the
difference of two group-algebra elements shows that the definitions
respect every linear relation. The maps are therefore well-defined
linear isometries. They extend to contractions $A,B:E\to E$, and
$(A+B)/2=I_E$ because $y+z=2x$. Lemma~\ref{lem:identity} gives
$A=B=I_E$, whence $y=z=x$.

The final assertion follows by first making the representation isometric
and then applying Theorem~\ref{thm:criterion}.
\end{proof}

\begin{corollary}[Finite-index extension]\label{cor:virtual}
Every uniformly bounded representation of a virtually FC group on a
separable real Banach space admits an equivalent invariant strictly convex norm.
\end{corollary}
\begin{proof}
This is the finite-index case of \cite[Lemma~3.2 and
Proposition~5.1]{DD}; we give a direct proof by averaging.
Let $H\leq G$ be FC with finite index $d$, and first make the given
representation isometric using $\|\cdot\|_0$.
Theorem~\ref{thm:fc} gives an equivalent $H$-invariant strictly convex norm $N_H$.
Choose representatives $r_1,\ldots,r_d$ for the right cosets
$H\backslash G$ and set
\[
 N_G(v)=\frac1d\sum_{j=1}^d N_H(\pi(r_j)v).
\]
Right multiplication permutes the right cosets, so the
$H$-invariance of $N_H$ makes $N_G$ invariant under $G$.
Each summand is an equivalent strictly convex norm, and their finite
average has the same properties. Normality of $H$ is not required.
\end{proof}

\section{The Dantas--Doucha obstruction}\label{sec:obstruction}
We first explain the analytic target of the converse. It is enough to
find $f$ and two sequences of probability measures such that
\begin{equation}\label{eq:target}
 f*\nu_n^*\longrightarrow f+\delta_e,\qquad
 f*(\nu_n')^*\longrightarrow f-\delta_e
\end{equation}
in supremum norm, with summable errors. Their average makes $f$
approximately harmonic, but their difference has a nonzero limit.
Equally, \eqref{eq:target} places $f$ at a nontrivial midpoint of its
closed convex orbit hull. This is the construction of Dantas--Doucha
in \cite[Example~5.5, Claim~2]{DD}.

\subsection{The three product properties}
We retain the notation of \cite[Example~5.5]{DD}. Let $\mathcal F$
be the set of finite sequences of distinct positive integers, including
the empty sequence $\varnothing$. Write $|v|$ for the length of $v$,
$v(i)$ for its $i$th entry, and $\ran(v)$ for its set of entries.
Given $(s_i)_{i\geq1}\subseteq G$, put
\[
 s_v=s_{v(1)}\cdots s_{v(|v|)},\qquad
 s_\varnothing=e,\qquad \mathcal S=\{s_v:v\in\mathcal F\}.
\]
Concatenation is denoted by juxtaposition; it belongs to $\mathcal F$
only when no entry is repeated. Equalities of sequences mean equality
as ordered strings, while equalities of their evaluations hold in $G$.

The hypotheses below are exactly the three properties in
\cite[Example~5.5, Claim~1]{DD}. Their roles can be seen before any
estimate. Injectivity makes the first entry of a nonempty product
well-defined. Appending an unused entry preserves that first entry;
appending an entry already used leaves $\mathcal S$. Finally, a point
outside $\mathcal S$ has at most one right successor in $\mathcal S$,
so averages over large blocks are small there. Thus a sign assigned to
the first entry is almost preserved by right averaging, except at the
empty sequence, where the first sign is created.

The needed approximation is uniform over all products. A fixed product
uses finitely many entries, but products of arbitrarily large length
can exhaust an entire averaging block. Give each letter a small
positive cost and let $w$ be the total cost of a word. We damp its
first-letter sign by $1/(1+w)$. A word that has used many letters from
one block has a large cost, and therefore a small value. This is the
uniformity mechanism in the proof below.

\begin{proposition}[{Dantas--Doucha obstruction, extracted from
\cite[Example~5.5, Claim~2]{DD}}]\label{thm:certificate}
Let $G$ be countable and $(s_i)_{i\geq1}\subseteq G$. Assume that:
\begin{enumerate}
\item the map $v\mapsto s_v$ on $\mathcal F$ is injective;
\item if $i\in\ran(v)$, then $s_vs_i\notin\mathcal S$;
\item if $g\notin\mathcal S$, at most one index $i$ satisfies
$gs_i\in\mathcal S$.
\end{enumerate}
There are $f\in\ell^\infty(G)$ and
$\mu_n,\nu_n,\nu_n'\in\Prob(G)$ with $\mu_n=(\nu_n+\nu_n')/2$ such that
\begin{equation}\label{eq:endpoint-estimates}
 \begin{aligned}
 \|f*\nu_n^*-(f+\delta_e)\|_\infty&\leq3\cdot2^{-n},\\
 \|f*(\nu_n')^*-(f-\delta_e)\|_\infty&\leq3\cdot2^{-n}.
 \end{aligned}
\end{equation}
Consequently $G$ is not positively strongly Choquet--Deny. The
right-translation representation on the separable cyclic space
\[
 X_f=\overline{\Span}\{R_gf:g\in G\}\subseteq\ell^\infty(G)
\]
has no equivalent invariant strictly convex norm; more precisely,
$f+\delta_e,f-\delta_e\in\overline{\co}(\Orb_G(f))$.
\end{proposition}
\begin{proof}
We adapt the argument of \cite[Example~5.5, Claim~2]{DD}. We retain
their blocks, signs, and additive weight, but use reciprocal damping
in place of their truncated affine factor. Blocks of size $2^{2n}$
suffice in place of $2^{2^n}$.
Partition $\N$ into consecutive blocks $N_n$ with $|N_n|=2^{2n}$,
and put $S_n=\{s_i:i\in N_n\}$. Split $S_n$ into two equally sized
sets $S_n^+,S_n^-$. Let $\mu_n,\nu_n,\nu_n'$ be uniform on
$S_n,S_n^+,S_n^-$, respectively. These are probability measures and
$\mu_n=(\nu_n+\nu_n')/2$.

Set $\sigma(s_i)=1$ on $S_n^+$ and $\sigma(s_i)=-1$ on $S_n^-$.
Following the notation of \cite{DD}, define the additive weight
$f':\mathcal S\to[0,\infty)$ by
\[
 f'(s_v)=\sum_{k\geq1}2^{-k}|N_k\cap\ran(v)|.
\]
The sum has finitely many nonzero terms, and the first hypothesis
makes it well-defined. Define
\[
 f(s_v)=\frac{\sigma(s_{v(1)})}{1+f'(s_v)}
   \quad(v\ne\varnothing),\qquad
 f(e)=0,\qquad
 f(g)=0\quad(g\notin\mathcal S).
\]
Then $\|f\|_\infty=1$, since $|f(s_i)|=(1+2^{-n})^{-1}$ for $s_i\in S_n$.
Fix $n$, and put $\varepsilon=2^{-n}$ and
$L=|S_n^+|=|S_n^-|=2^{2n-1}$. We estimate the two averages in
\eqref{eq:endpoint-estimates} by the same argument.
At $e$ their values are $(1+\varepsilon)^{-1}$ and
$-(1+\varepsilon)^{-1}$, so each has error at most $\varepsilon$
from its target. Outside $\mathcal S$,
at most one successor contributes, and the error is at most
$L^{-1}=2\varepsilon^2\leq\varepsilon$.

Let $g=s_v\ne e$, write $w=f'(s_v)$, and consider either half-block.
Let $r$ be the number of its entries already used in $v$. Then
$r\varepsilon\leq w$. Each used entry gives a zero successor by the
second hypothesis. Every unused entry preserves the first letter and
adds $\varepsilon$ to the weight. Writing
$\sigma=\sigma(s_{v(1)})$, the half-block average is therefore
\[
 \frac{(1-r/L)\sigma}{1+w+\varepsilon}.
\]
Its error from $f(g)=\sigma/(1+w)$ is exactly
\[
 \frac{\varepsilon}{(1+w)(1+w+\varepsilon)}
 +\frac{r}{L(1+w+\varepsilon)}
 \leq\varepsilon+\frac1{L\varepsilon}=3\varepsilon,
\]
where we used $r\varepsilon\leq w\leq1+w+\varepsilon$.
This single estimate covers all word lengths and both half-blocks,
and proves \eqref{eq:endpoint-estimates}.

Taking the average and difference now gives
\begin{align}
 \|f*\mu_n^*-f\|_\infty&\leq3\cdot2^{-n},\label{eq:main-approx}\\
 \|d_n-2\delta_e\|_\infty&\leq6\cdot2^{-n},
 \qquad d_n=f*(\nu_n^*-(\nu_n')^*).\label{eq:main-split}
\end{align}
In particular,
\[
 \sum_{n\geq1}\|d_{n+1}-d_n\|_\infty
 \leq6\sum_{n\geq1}(2^{-(n+1)}+2^{-n})<\infty,
 \qquad d_n\longrightarrow2\delta_e\ne0.
\]
This disproves the positive strong condition.

The endpoint estimates also give the geometric conclusion directly.
The space $X_f$ is separable because $G$ is countable, and is invariant
under the isometries $R_g$. Each probability average is a convex
combination of right translates of $f$. Hence \eqref{eq:endpoint-estimates}
places the two distinct limits $f+\delta_e,f-\delta_e$ in
$\overline{\co}(\Orb_G(f))$. Their midpoint is $f$, which precludes
an invariant strictly convex norm.
\end{proof}

\section{Realizing the product conditions}\label{sec:sequence}
It remains to produce a sequence satisfying the three hypotheses of
Proposition~\ref{thm:certificate} when $[G:\FC(G)]=\infty$.
Throughout this section put $F=\FC(G)$. The group need not be countable.

\subsection{Rewritability and the product conditions}
Finite index of the FC-center forces collisions among products.
The following observation uses the prefix-coset argument in
\cite[Lemmas~3.1--3.2]{BR}.

\begin{lemma}\label{lem:block-commutation}
If $[G:\FC(G)]<\infty$, every infinite sequence in $G$ has two
disjoint nonempty finite consecutive blocks whose products commute.
In particular, its products of distinct-index words cannot all be
distinct.
\end{lemma}
\begin{proof}
Write the sequence as $(s_i)$ and its prefix products as
$p_k=s_1\cdots s_k$, with $p_0=e$. Two prefixes lie in the same
coset of $F$, say $p_aF=p_bF$ with $a<b$. Hence the nonempty block
product
\[
 u=p_a^{-1}p_b=s_{a+1}\cdots s_b
\]
belongs to $F$. Now $C_G(u)$ has finite index. Apply the same
pigeonhole argument to prefixes $q_k=s_{b+1}\cdots s_{b+k}$ of the
remaining tail. Equality $q_cC_G(u)=q_dC_G(u)$, with $c<d$, gives
\[
 v=q_c^{-1}q_d=s_{b+c+1}\cdots s_{b+d}\in C_G(u).
\]
Thus $uv=vu$, while the two corresponding concatenations of the
disjoint blocks are different words with no repeated index.
\end{proof}

Blyth--Rhemtulla call a group \emph{eventually rewritable}, or a
$\mathcal Q_\infty$-group, if every sequence $(x_i)_{i\geq1}$ admits $n\geq2$
and distinct permutations $\sigma,\tau$ of $\{1,\ldots,n\}$ such that
\[
 x_{\sigma(1)}\cdots x_{\sigma(n)}
 =x_{\tau(1)}\cdots x_{\tau(n)}.
\]
They prove that this is equivalent to being FC-by-finite, that is,
virtually FC \cite[Proposition~2]{BR}. Consequently, outside this
class there is a sequence whose products distinguish every ordering
of each finite initial segment. They also distinguish orderings of
any fixed finite subset: append its complement in an initial segment
and cancel. This explains why products that retain their order are
natural here. It does not yet distinguish different supports, exclude
repeated appends, or control right successors outside the product set.
Those are the extra requirements of Proposition~\ref{thm:certificate}.

Neither extra requirement follows merely from distinguishing ordered
words. The standard generators of the free product
$\ast_{i\geq1}\Z_2$ distinguish all words with distinct indices, but
$s_i^2=e$ violates the repeated-append condition. Passing to a
subsequence does not fix this.

The exterior-successor condition is also independent of unique decoding.
For example, in the free group on $a,b,t,c_1,c_2,\ldots$, the sequence
\[
 a,\ b,\ ta,\ tb,\ c_1,\ c_2,\ldots
\]
has distinct values for all positive words, even when indices repeat:
read a symbol $t$ together with the following $a$ or $b$, and every
other symbol separately. This uniquely decodes a product into the
displayed generators, and also gives the repeated-append exclusion.
Nevertheless, $t$ is outside their positive
product set and has the two successors $ta,tb$ in it. Thus even
unique decoding of all positive words would not by itself supply
the analytic obstruction.

The basic induction is also suggested by the proof of
\cite[Lemma~3.3]{BR}. If a new element $x$ is inserted into two old
products, a collision takes the form
\[
 axb=cxd\quad\Longleftrightarrow\quad
 x(bd^{-1})x^{-1}=a^{-1}c.
\]
When $bd^{-1}\notin F$, this excludes at most one coset of an
infinite-index centralizer. When $bd^{-1}\in F$, its centralizer has
finite index; we instead choose $x$ in that centralizer. The equation
then reduces to an equality among old products. The induction must
retain enough information to prove that each such surviving equality
is harmless. The finite extension lemma below uses two necessary
features of an infinite sequence with the desired predecessor
property: nonempty products avoid $F$, and different letters lie in
different cosets of $F$.

This use of centralizer cosets is shared with the switching literature.
The ordinary switching argument is explicit in
\cite[Claim~3.3]{FTVF}; super-switching adds equations of the form
$xax=b$ in \cite[Definition~2.3 and Proposition~2.5]{FHTVF}.
Erschler--Kaimanovich organize these arguments using negligible sets
and prove super-switching existence for all ICC groups
\cite[Sections~4.B--4.C, especially Proposition~4.25]{EK}.
We use their notation below. Our group can have a nontrivial FC-center,
and the common-predecessor requirement also produces
$xax^{-1}=bxc$. We therefore prove the precise simultaneous avoidance
statement needed here, rather than assume it follows from a switching
statement for ICC groups.

\subsection{Negligible sets and finite equation avoidance}
Following \cite[Definition~4.9]{EK}, a subset of $G$ is
\emph{negligible} if it is contained in a finite union of cosets of
infinite-index subgroups. This family is closed under subsets, finite
unions, and left and right translation. For right translation, use
$aKb=ab(b^{-1}Kb)$.

\begin{lemma}[{Neumann's finite coset covering lemma \cite{Neumann}}]
\label{lem:coset-cover}
No finite-index subgroup of $G$ is negligible.
\end{lemma}
\begin{proof}
Neumann's theorem says that a group covered by finitely many cosets of
subgroups has a finite-index subgroup among them; see also
\cite[Lemma~4.13]{EK}. For the stated version, let $H\leq G$ have
finite index. Each nonempty $H\cap aK$ is a coset in $H$ of $H\cap K$.
If $[G:K]=\infty$, then $[H:H\cap K]=\infty$. A negligible cover of
$H$ would therefore contradict the theorem applied to $H$.
\end{proof}

Following \cite[\S\S4.B--4.C]{EK}, write the conjugacy and quadratic
solution sets as
\[
 S_{a\to b}=\{x\in G:xax^{-1}=b\},\qquad
 S^{\pm}_{a\to b}=\{x\in G:xax=b\}.
\]
If $x_0\in S_{a\to b}$, then
\begin{equation}\label{eq:solution-coset}
 S_{a\to b}=x_0C_G(a).
\end{equation}
Indeed, $xax^{-1}=x_0ax_0^{-1}$ if and only if $x_0^{-1}x$ commutes
with $a$. Thus this solution set is negligible whenever $a\notin F$.
Erschler--Kaimanovich also show that $S^{\pm}_{a\to b}$ is contained
in a coset of $C_G(ab)$ \cite[Lemma~4.23]{EK}. This handles $ab\notin F$,
but our induction must allow arbitrary coefficients, including the
involution equation $x^2=e$. The following overlap argument treats
all the required quadratic and mixed equations at once.

\begin{lemma}[Finite equation avoidance]\label{lem:avoid}
Suppose $[G:F]=\infty$ and $H\leq G$ has finite index. No finite union
of a negligible set and sets of the following two forms covers $H$:
\[
 S^{\pm}_{a\to b}=\{x:xax=b\},\qquad
 \{x:xax^{-1}=bxc\},\qquad a,b,c\in G.
\]
\end{lemma}
\begin{proof}
We first show that if $E$ is either displayed solution set and
$h\notin F$, then $E\cap h^{-1}E$ is negligible. Membership means
$x,hx\in E$. In the quadratic case this gives
\[
 hxahx=xax
 \quad\Longrightarrow\quad
 hxah=xa
 \quad\Longrightarrow\quad
 x(aha^{-1})x^{-1}=h^{-1}.
\]
Thus $E\cap h^{-1}E\subseteq S_{aha^{-1}\to h^{-1}}$.
The element $aha^{-1}$ is outside $F$, so
\eqref{eq:solution-coset} puts the intersection in a coset of an
infinite-index centralizer, unless it is empty.

For the mixed equation $xax^{-1}=bxc$, substituting this relation
into the equation for $hx$ yields
\[
 hbxch^{-1}=bhxc
 \quad\Longrightarrow\quad
 x(ch^{-1}c^{-1})x^{-1}=b^{-1}h^{-1}bh.
\]
Since $ch^{-1}c^{-1}\notin F$, the same reasoning proves negligibility
of the intersection. This is a statement about the overlaps; it does
not assert that $E$ itself is negligible.

Suppose now that $H\subseteq A\cup E_1\cup\cdots\cup E_r$, where
$A$ is negligible and the $E_i$ are solution sets of the displayed
forms. Since $[H:H\cap F]=\infty$, choose $h_0,\ldots,h_r\in H$
in distinct cosets modulo $F$. For each $x\in H$, either one of
$h_jx$ belongs to $A$, or two of the $r+1$ points belong to the same
$E_i$. Hence
\[
 H\subseteq\bigcup_{j=0}^r h_j^{-1}A\ \cup\!
 \bigcup_{\substack{1\leq i\leq r\\0\leq j<k\leq r}}
 (h_j^{-1}E_i\cap h_k^{-1}E_i).
\]
The pair intersection equals
$h_j^{-1}(E_i\cap(h_kh_j^{-1})^{-1}E_i)$ and is negligible by the
first part, since $h_kh_j^{-1}\notin F$. The right-hand side is thus
negligible, contradicting Lemma~\ref{lem:coset-cover}.
\end{proof}

Related recursive avoidance arguments occur in
\cite[Theorems~4 and~5]{BP}, using centralizers and group Zariski
topologies. No topological nondiscreteness assumption is needed in
Lemma~\ref{lem:avoid}.

\subsection{The finite extension problem}
The three product requirements follow from a single condition,
if we allow the identity as an empty terminal. Set $s_0=e$ and use the
convention $t0=t$. The symbol $0$ is never an entry of a word in
$\mathcal F$; it is used only for the terminal indices below.

For a finite stage, let $\mathcal F_n$ be the words with entries in
$\{1,\ldots,n\}$ and $\mathcal S_n=\{s_v:v\in\mathcal F_n\}$.
Consider the following \emph{common-predecessor condition}:
\begin{equation}\label{eq:predecessor}
 \begin{gathered}
 v,w\in\mathcal F_n,\quad i,j\in\{0,\ldots,n\},\quad i\ne j,
 \quad s_vs_i^{-1}=s_ws_j^{-1}\\
 \Longrightarrow\quad v=ti,\quad w=tj
 \quad\text{for some }t\in\mathcal F_n.
 \end{gathered}
\end{equation}
In words, if a point has right successors in the product set indexed
by two distinct terminals, both must come from appending those terminals to one
word. Allowing terminal $0$ also handles a point already in the
product set.

Condition~\eqref{eq:predecessor} is an exact reformulation of the
three product requirements, valid also at a finite stage. First,
it implies each of them:
\begin{enumerate}
\item If $s_v=s_w$ and $v=ti$ is nonempty, then
$s_ws_i^{-1}=s_ts_0^{-1}$. The condition gives $w=ui$ and $t=u$,
so $w=v$. Interchanging $v,w$ if necessary proves injectivity;
the case of two empty words is immediate.
\item If $s_vs_i=s_w$ with $1\leq i\leq n$, then
$s_ws_i^{-1}=s_vs_0^{-1}$. The condition gives $w=vi$ as a word
without repeated entries. Hence $i\notin\ran(v)$.
\item If $gs_i=s_v$ and $gs_j=s_w$ with distinct $i,j\in\{1,\ldots,n\}$, then the condition
gives $v=ti$, $w=tj$, and therefore $g=s_t\in\mathcal S_n$.
\end{enumerate}
Conversely, suppose the three requirements hold and put
$g=s_vs_i^{-1}=s_ws_j^{-1}$. If one terminal is $0$, then
$g\in\mathcal S_n$ immediately. Otherwise, the exterior-successor
condition gives the same conclusion. Write $g=s_t$.
The repeated-append exclusion says that each nonzero terminal is
absent from $t$, and injectivity now gives $v=ti$, $w=tj$.
The same reasoning applies to an infinite sequence.

It remains to extend finite tuples while preserving
\eqref{eq:predecessor}. Two further conditions will be maintained:
\begin{align}
 s_v&\notin F\quad
       (v\in\mathcal F_n\setminus\{\varnothing\}),
       \label{eq:fc-avoidance}\\
 s_iF&\ne s_jF\quad(1\leq i<j\leq n).
       \label{eq:coset-separation}
\end{align}
Both conditions are necessary for any infinite sequence satisfying
\eqref{eq:predecessor}. For \eqref{eq:fc-avoidance}, suppose a nonempty
word $U$ has value $u\in F$. A nonempty block $T$ using only later
indices can be chosen with $s_T\in C_G(u)$, by the tail-prefix
pigeonhole argument of Lemma~\ref{lem:block-commutation}.
Then $s_{UT}=s_{TU}$ contradicts injectivity.
For \eqref{eq:coset-separation}, if $h=s_i^{-1}s_j\in F$, choose
such a later block $T$ with $s_T\in C_G(h)$. Then
\[
 s_{jT}s_j^{-1}=s_j s_Ts_j^{-1}
 =s_i s_Ts_i^{-1}=s_{iT}s_i^{-1}.
\]
Condition~\eqref{eq:predecessor} would require the word $jT$ to end
in $j$, which it does not.

For a finite tuple these are additional hypotheses, since the tail
used above is not available. In the extension proof they distinguish
the possible insertion positions of the new letter.

\begin{lemma}[Finite extension]\label{lem:extension}
Suppose $[G:F]=\infty$, $n\geq0$, and a finite tuple
$s_1,\ldots,s_n$ satisfies \eqref{eq:predecessor},
\eqref{eq:fc-avoidance}, and \eqref{eq:coset-separation}.
There is an element $x\in G$ such that the augmented tuple
$s_1,\ldots,s_n,x$ satisfies all three conditions.
\end{lemma}
\begin{proof}
Write $\star$ for the new formal letter $n+1$ and $x$ for its value.
A new positive word has the form $A\star B$, where $A,B$ are old
words with disjoint sets of entries. Lowercase $a,b,c,d$ denote the
evaluations of old words $A,B,C,D$.

\smallskip\noindent
\emph{How conjugacy equations are treated.}
For each conjugacy equation
\begin{equation}\label{eq:conjugacy}
 xux^{-1}=z
\end{equation}
encountered below, if $u\notin F$ we exclude its solution set, which
is empty or a coset of the infinite-index centralizer $C_G(u)$.
For the finitely many coefficients $u\in F$, require $x\in C_G(u)$.
Their centralizers have a finite-index intersection $H$; take $H=G$
if there are none. All coefficients are determined by the old
sequence, so $H$ is fixed before $x$ is chosen. Within $H$, an equation
\eqref{eq:conjugacy} with $u\in F$ becomes the old equality $u=z$:
the new letter has disappeared. We check below that these surviving
equalities give precisely the required predecessor.
All other forbidden sets are singletons, cosets of $F$, or the
quadratic and mixed solution sets in Lemma~\ref{lem:avoid}.

\smallskip\noindent
\emph{Preserving the two FC-center conditions.}
A new positive product has the form $axb$, so
\eqref{eq:fc-avoidance} excludes the sets
\[
 \{x:axb\in F\}=a^{-1}Fb^{-1}.
\]
Condition~\eqref{eq:coset-separation} also excludes $s_1F,\ldots,s_nF$.
Since $F$ is normal and has infinite index, all of these are
infinite-index cosets, and there are only finitely many of them.

\smallskip\noindent
\emph{Both distinguished terminals old.}
Take distinct $i,j\in\{0,\ldots,n\}$ and consider
$s_vs_i^{-1}=s_ws_j^{-1}$. If neither word contains $\star$, the old
condition applies. If exactly one does, the equation fixes a single
value of $x$. If both do, write $v=A\star B$, $w=C\star D$; then
\[
 axb s_i^{-1}=cxd s_j^{-1}
 \quad\Longleftrightarrow\quad
 x(bs_i^{-1}s_jd^{-1})x^{-1}=a^{-1}c.
\]
The choice rule leaves only
$bs_i^{-1}s_jd^{-1}=a^{-1}c\in F$. Thus
$(ab)s_i^{-1}=(cd)s_j^{-1}$, and the old predecessor condition gives
$AB=Ti$ and $CD=Tj$ for an old word $T$.

Now $A,C$ are prefixes of $Ti,Tj$, and $a^{-1}c\in F$.
If one is a proper prefix of the other, $a^{-1}c$ is the value of
a nonempty old word or the inverse of such a value, contradicting
\eqref{eq:fc-avoidance}. If they are incomparable, they must be
$A=Ti$, $C=Tj$, with $i,j\ne0$. In this case
$a^{-1}c=s_i^{-1}s_j\notin F$ by \eqref{eq:coset-separation},
again a contradiction. Hence $A=C$, and this common prefix lies
in $T$. Inserting $\star$ at that position in $T$ gives
$\widetilde T$ with $v=\widetilde Ti$ and $w=\widetilde Tj$,
as required. This also covers either terminal being $0$.

\smallskip\noindent
\emph{One distinguished terminal new.}
By symmetry it remains to consider
\begin{equation}\label{eq:newterminal}
 s_vx^{-1}=s_ws_j^{-1},\qquad j\in\{0,\ldots,n\}.
\end{equation}
There are four possibilities for the occurrence of $\star$ in $v,w$.
If neither word contains it, the equation fixes $x$.
If only $v=A\star B$ contains it and $w=C$, the equation is
\[
 xbx^{-1}=a^{-1}c s_j^{-1}.
\]
If $B$ is nonempty, \eqref{eq:fc-avoidance} gives $b\notin F$,
so this solution set is excluded. If $B$ is empty, the equation is
$as_j=c$. For $j\ne0$, the old predecessor condition with terminals
$0,j$ gives $C=Aj$. For $j=0$, old injectivity, already deduced from
that condition, gives $C=A$. Thus in either case
$v=A\star$, $w=Aj$, exactly the required pattern.

If only $w=C\star D$ contains $\star$ and $v=A$, then
\[
 ax^{-1}=cxd s_j^{-1}
 \quad\Longleftrightarrow\quad x(ds_j^{-1})x=c^{-1}a,
\]
which is quadratic. If both words contain $\star$, then
\[
 axbx^{-1}=cxd s_j^{-1}
 \quad\Longleftrightarrow\quad
 xbx^{-1}=a^{-1}c\,x\,(ds_j^{-1}),
\]
which is mixed. Lemma~\ref{lem:avoid} allows us to exclude both types.

Every possible failure of the three maintained conditions has now been
accounted for. There are only finitely many equations at each stage,
and singletons are negligible because $G$ is infinite. Applying
Lemma~\ref{lem:avoid} within $H$ supplies the next element $x$.
This proves the finite extension statement.
\end{proof}

\begin{theorem}\label{thm:sequence}
If $[G:\FC(G)]=\infty$, there is a sequence $(s_i)_{i\geq1}$
satisfying the three product hypotheses of
Proposition~\ref{thm:certificate}. No countability assumption on $G$
is needed for this construction.
\end{theorem}
\begin{proof}
The empty tuple satisfies \eqref{eq:predecessor},
\eqref{eq:fc-avoidance}, and \eqref{eq:coset-separation}
vacuously. Apply Lemma~\ref{lem:extension}
successively. Any failure of \eqref{eq:predecessor} for the resulting
infinite sequence would involve only finitely many entries and
terminal indices, so would already occur at a finite stage. The
three deductions above therefore give all the required product
properties.
\end{proof}

\section{The classification and examples}\label{sec:examples}
The preceding sections now give both implications of the main theorem.

\begin{proof}[Proof of Theorem~\ref{thm:classification}]
Conditions (2) and (3) are equivalent by
\cite[Theorem~4.2]{DD}. Corollary~\ref{cor:virtual} proves
(1)$\Rightarrow$(2), and restricting to probability measures gives
(3)$\Rightarrow$(4).

If (1) fails, \eqref{eq:fc-index} gives $[G:\FC(G)]=\infty$.
Theorem~\ref{thm:sequence} supplies the sequence required by
Proposition~\ref{thm:certificate}, whose probability measures violate
(4). This proves (4)$\Rightarrow$(1) and completes the equivalence.
\end{proof}

\subsection{Permanence and examples}
In particular, among countable groups the strongly Choquet--Deny groups
are exactly the eventually rewritable ($\mathcal Q_\infty$) groups, by
\cite[Proposition~2]{BR} and Theorem~\ref{thm:classification}.
The following consequences distinguish the classification from both
virtual abelianness and the ordinary Choquet--Deny property.

\begin{corollary}\label{cor:permanence}
Among countable groups, the strongly Choquet--Deny property is preserved
by subgroups, quotients, finite direct products, and passage to a group
containing the given group with finite index.
\end{corollary}
\begin{proof}
These are permanence properties of virtually FC groups. If $H\leq G$
is FC of finite index, then $K\cap H$ is FC of finite index in every
subgroup $K\leq G$, and the image of $H$ in any quotient is FC of
finite index. Products and finite-index overgroups are immediate.
Apply Theorem~\ref{thm:classification}.
\end{proof}

The subgroup and finite-index permanence were already established in
\cite[Proposition 5.1]{DD}. The classification also makes clear why
increasing unions need not preserve the property.

\begin{example}\label{ex:finite-sum}
Every countable direct sum of finite groups is strongly Choquet--Deny:
conjugation preserves the finite support of each element, so the group
is FC. More generally, a countable increasing union of finite normal
subgroups is FC, since the conjugacy class of each element stays in
one of those subgroups.

In particular, $G_0=\bigoplus_{n\geq1}S_3$ is strongly Choquet--Deny
but is not virtually abelian. To see the latter assertion, let
$F_n=\prod_{j=1}^n S_3$ and let $A\leq G_0$ be abelian. Each coordinate
projection of $A\cap F_n$ has order at most three, so
\[
 |A\cap F_n|\leq3^n,\qquad
 [G_0:A]\geq[F_n:A\cap F_n]\geq2^n.
\]
Thus $A$ has infinite index.
\end{example}

\begin{example}\label{ex:dihedral}
Let $D_\infty=\Z\rtimes(\Z/2\Z)$, where the nontrivial element acts
by inversion, and put $G=\bigoplus_{n\geq1}D_\infty$.
The FC-center of $D_\infty$ is its translation subgroup $\Z$:
a translation has at most two conjugates, whereas every reflection
has infinitely many. Conjugation in $G$ is coordinatewise, and hence
\[
 \FC(G)=\bigoplus_{n\geq1}\Z,\qquad
 G/\FC(G)\cong\bigoplus_{n\geq1}\Z/2\Z.
\]
The FC-center has infinite index, so $G$ is not strongly Choquet--Deny.
Nevertheless, $G$ is an increasing union of the normal direct factors
$\prod_{j=1}^nD_\infty$, each of which is virtually abelian and strongly
Choquet--Deny. Thus neither increasing unions of normal strongly
Choquet--Deny subgroups nor countable direct sums preserve the property.

This group is ordinary Choquet--Deny. Indeed, in an ICC quotient the
image of $\FC(G)$ must be trivial, because its elements retain finite
conjugacy classes. Such a quotient would then be abelian, hence trivial.
The assertion follows from \cite[Theorem 1]{FHTVF}.
\end{example}

\begin{example}\label{ex:heisenberg}
The integral Heisenberg obstruction is \cite[Example~5.5]{DD}.
The following locally finite variant illustrates why the full countable
classification requires more than the finitely generated case.
Let $p$ be a prime and $R=\F_p[t]$. The Heisenberg group
\[
 H_3(R)=\left\{\begin{pmatrix}1&a&c\\0&1&b\\0&0&1\end{pmatrix}:
 a,b,c\in R\right\}
\]
is countable, locally finite, and nilpotent of class two. For local
finiteness, the first and second coordinates of a finitely generated
subgroup lie in finite-dimensional $\F_p$-spaces, and its third
coordinates lie in the span of the original third coordinates and
the pairwise products of the first and second coordinates.

Its FC-center is exactly its center, consisting of the elements with
$a=b=0$. Indeed, if $a\ne0$, conjugation by the elements with coordinates
$(0,b',0)$ gives infinitely many distinct shifts $ab'$ in the central coordinate.
If $b\ne0$, the analogous conjugations by $(a',0,0)$ give the same
conclusion. We use here that $R$ is an infinite integral domain.
Thus $[H_3(R):\FC(H_3(R))]=\infty$, and the group is not strongly
Choquet--Deny.

Every nontrivial quotient is nilpotent of class at most two and has
nontrivial center, so no quotient is ICC. It follows again from
\cite[Theorem 1]{FHTVF} that $H_3(R)$ is ordinary Choquet--Deny.
This gives a locally finite example separating the two properties.
\end{example}

\begin{thebibliography}{FHTVF}
\bibitem[BP]{BP}
T.~Banakh and I.~Protasov,
\emph{Zariski topologies on groups},
preprint, 2010,
\href{https://arxiv.org/abs/1001.0601v1}{arXiv:1001.0601v1}.

\bibitem[BR]{BR}
R.~D.~Blyth and A.~H.~Rhemtulla,
\emph{Rewritable products in FC-by-finite groups},
Canad. J. Math. \textbf{41} (1989), no.~2, 369--384.
\href{https://doi.org/10.4153/CJM-1989-018-8}{doi:10.4153/CJM-1989-018-8}.

\bibitem[DD]{DD}
S.~Dantas and M.~Doucha,
\emph{Strongly Choquet--Deny groups and invariant strictly convex renormings},
preprint, 2026,
\href{https://arxiv.org/abs/2610.23584v1}{arXiv:2610.23584v1}.
All numbered citations refer to this version.

\bibitem[D]{D}
M.~Doucha,
\emph{Invariant strictly convex renormings},
J. Funct. Anal. \textbf{290} (2026), no.~11, article~111428.
\href{https://doi.org/10.1016/j.jfa.2026.111428}{doi:10.1016/j.jfa.2026.111428}.

\bibitem[EFHN]{EFHN}
T.~Eisner, B.~Farkas, M.~Haase, and R.~Nagel,
\emph{Operator Theoretic Aspects of Ergodic Theory},
Graduate Texts in Mathematics, vol.~272, Springer, Cham, 2015.
\href{https://doi.org/10.1007/978-3-319-16898-2}{doi:10.1007/978-3-319-16898-2}.

\bibitem[EK]{EK}
A.~Erschler and V.~A.~Kaimanovich,
\emph{Arboreal structures on groups and the associated boundaries},
Geom. Funct. Anal. \textbf{33} (2023), 694--748.
\href{https://doi.org/10.1007/s00039-023-00633-y}{doi:10.1007/s00039-023-00633-y}.
Section, equation, and result numbers refer to the accessible preprint
\href{https://arxiv.org/abs/1903.02095v1}{arXiv:1903.02095v1} (2019).

\bibitem[FHTVF]{FHTVF}
J.~Frisch, Y.~Hartman, O.~Tamuz, and P.~Vahidi Ferdowsi,
\emph{Choquet--Deny groups and the infinite conjugacy class property},
Ann. of Math. (2) \textbf{190} (2019), no.~1, 307--320.
\href{https://doi.org/10.4007/annals.2019.190.1.5}{doi:10.4007/annals.2019.190.1.5}.

\bibitem[FTVF]{FTVF}
J.~Frisch, O.~Tamuz, and P.~Vahidi Ferdowsi,
\emph{Strong amenability and the infinite conjugacy class property},
Invent. Math. \textbf{218} (2019), 833--851.
\href{https://doi.org/10.1007/s00222-019-00896-z}{doi:10.1007/s00222-019-00896-z}.
\href{https://arxiv.org/abs/1801.04024v5}{arXiv:1801.04024v5}.

\bibitem[J]{J}
W.~Jaworski,
\emph{Countable amenable identity excluding groups},
Canad. Math. Bull. \textbf{47} (2004), no.~2, 215--228.
\href{https://doi.org/10.4153/CMB-2004-021-1}{doi:10.4153/CMB-2004-021-1}.

\bibitem[JN]{JN}
W.~Jaworski and M.~Neufang,
\emph{The Choquet--Deny equation in a Banach space},
Canad. J. Math. \textbf{59} (2007), no.~4, 795--827.
\href{https://doi.org/10.4153/CJM-2007-034-4}{doi:10.4153/CJM-2007-034-4}.

\bibitem[Neu]{Neumann}
B.~H.~Neumann,
\emph{Groups covered by finitely many cosets},
Publ. Math. Debrecen \textbf{3} (1954), 227--242.
\href{https://doi.org/10.5486/PMD.1954.3.3-4.07}{doi:10.5486/PMD.1954.3.3-4.07}.
\end{thebibliography}
\end{document}
